\documentclass[10pt]{article}
\usepackage{graphicx}
\setlength{\textwidth}{18cm}
\setlength{\textheight}{25cm}
\setlength{\oddsidemargin}{-1cm}
\setlength{\topmargin}{-1.5cm}
\newcommand{\farcs}{\mathord{.\!^{\prime\prime}}}
\pagestyle{empty}
\begin{document}
\setcounter{figure}{7}
\begin{figure}[p]
    \centering
    \includegraphics[width=\textwidth]{figs/shape_dependence_v04.pdf}
    \caption{Observed-shape recovery versus input NISP weighted S/N.
    Left and middle: percentage reduction in ellipticity-component RMSE relative to native NISP, for fixed weighted and adaptive moments, respectively; positive values indicate lower error against observed JWST F115W.
    Ochre circles, cobalt squares, and green triangles denote a single diffusion realization, ResNet, and ResNet-D5.
    Each estimator/bin uses the same galaxies measurable in JWST, native NISP, and all six reconstruction/interpolation methods.
    Points are horizontally offset for visibility; the bins are categorical and connecting lines guide the eye.
    Whiskers show conditional 95\% intervals from 2000 paired spatial-block resamples, not simultaneous intervals for a trend.
    Right: common measurement coverage among all attempted fields, with fixed moments in solid orange and adaptive moments in hatched charcoal.
    Numbers above bars are common successes; $N$ below each bin gives the attempted count.
    Totals are 1446 fixed-moment and 1498 adaptive-moment successes out of 2118 validation fields.
    The S/N proxy neglects correlated noise, and fixed and adaptive moments have different successful samples.
    These observed cross-band comparisons include PSF and pixel-sampling differences and do not measure calibrated shear.}
    \label{fig:shape_snr}
\end{figure}
\clearpage
\begin{figure}[p]
    \centering
    \includegraphics[width=\textwidth]{figs/shear_response_v06.pdf}
    \caption{Controlled shear response of observed adaptive shapes using the four supplied NISP effective PSFs.
    Panels (a) and (b) show exponential galaxies with intrinsic half-light radii $0\farcs18$ and $0\farcs35$.
    Charcoal circles, cobalt squares, and orange open diamonds denote native NISP, the original ResNet, and a noiseless high-resolution control.
    The vertical axis is $\mathrm{tr}(R)/\mathrm{tr}(R_{\rm intrinsic})$; the dotted line marks the intrinsic response.
    Each point equally averages the four per-PSF ratios on identical objects successful across all PSFs, methods, and five shear branches.
    Each noisy cell attempts eight orientations and four noise realizations per orientation; the common sample is 31/32 at radius $0\farcs18$ and S/N 20, and 32/32 elsewhere in the displayed grid.
    The same noise is reused across shear branches and PSFs.
    S/N positions are categorical and connecting lines guide the eye; these are point estimates without uncertainty intervals.
    The high-resolution control uses a Gaussian PSF of FWHM $0\farcs08$, not a validated JWST PSF.
    The two radii were selected for their larger gains in the earlier Gaussian-PSF pilot; the full new grid also includes radius $0\farcs65$ and noiseless native inputs.
    No PSF correction is applied to these shapes, and the responses do not establish calibrated weak-lensing accuracy.}
    \label{fig:shear_response}
\end{figure}
\clearpage
\end{document}
